\sectioncentered*{Заключение}
\addcontentsline{toc}{section}{Заключение}

В ходе выполнения курсового проекта была разработана база знаний для интеллектуального помощника, предназначенного для подготовки к техническим собеседованиям в IT-области.

\textbf{Выполнен анализ} предметной области подготовки к техническим интервью, в результате которого были выделены три основные области знаний: Computer Science, архитектура систем и процесс прохождения собеседования. Также были определены ключевые типы решаемых задач: диагностика уровня компетенций, классификация учебных материалов, построение учебных траекторий и генерация вопросов.

\textbf{Обоснован выбор} документно-ориентированного подхода на основе Markdown и YAML-фронтматтера, который обеспечивает сочетание удобочитаемости для человека, машиночитаемости и совместимости с системами контроля версий.

\textbf{Спроектированы} логическая и физическая модели базы знаний, включающие доменное разделение (\texttt{cs/}, \texttt{arch/}, \texttt{process/}, \texttt{global/}), систему уникальной идентификации и механизм семантических связей между единицами знаний.

\textbf{Создан} CLI-инструмент на Python 3.12 для управления базой знаний, поддерживающий поиск, валидацию метаданных, генерацию реестра и создание новых единиц знаний.

\textbf{Разработан} веб-MVP на базе FastAPI, предназначенный для интерактивной демонстрации работы базы знаний, включая поиск и навигацию по семантическим связям.

\textbf{Сформировано наполнение} базы знаний: на момент завершения проекта она включает 120 единиц знаний (20 по направлению CS, 20 по arch и 20 по process).

\textbf{Проведено комплексное тестирование}: подготовлено 177 автоматизированных тестов на базе \texttt{pytest}, охватывающих модули парсинга, валидации, поиска и генерации. Доля успешно пройденных тестов составляет около $\sim$98\%.

\textbf{Проведена оценка качества} базы знаний: время поиска по метаданным не превышает 2 секунд при выборке до 50 единиц, а в документах со статусом утверждённых отсутствуют критические ошибки валидации.

Таким образом, все задачи, сформулированные во введении, выполнены в полном объёме.

\textbf{Ожидаемый практический эффект:}
\begin{itemize}
    \item персонализация подготовки кандидатов за счёт выявления пробелов в знаниях;
    \item уменьшение временных затрат на поиск релевантных материалов;
    \item повышение общей эффективности подготовки к техническим собеседованиям.
\end{itemize}

\textbf{Перспективные направления развития:}
\begin{itemize}
    \item подключение RAG-конвейера (LangChain + ChromaDB) для реализации семантического поиска;
    \item создание модуля решателя задач для автоматизированного построения учебных траекторий;
    \item расширение базы знаний до целевого объёма $\geq$ 553 единиц.
\end{itemize}


\textbf{Индивидуальный вклад автора:}
\begin{itemize}
    \item анализ предметной области и существующих решений (HackerRank, Computer Science Ontology (CSO Portal), VisuAlgo);
    \item проектирование системы метаданных с 11 обязательными полями и таксономии тегов;
    \item реализация модулей \texttt{search.py}, \texttt{registry.py};
    \item разработка веб-MVP на FastAPI с эндпоинтом \texttt{/api/ask};
    \item создание части тестового набора (58 тестов pytest) и проведение тестирования;
    \item наполнение предметной области «Computer Science»:20 единиц знаний.
\end{itemize}

Разработанная база знаний выступает информационным ядром для следующих этапов курсового проектирования: предоставляет структурированные данные для модуля решателя задач, обеспечивает API для пользовательских интерфейсов и формирует датасет для обучения рекомендательных моделей.